\chapter{Testare și evaluare}

După definirea cerințelor și prezentarea implementării, următorul pas firesc în evaluarea soluției WhereWeFishin este verificarea comportamentului acesteia prin teste și prin validări operaționale. În contextul unui sistem care combină autentificare și autorizare, rezervări, administrare multi-rol, integrare cu servicii externe și procesare media, testarea nu poate fi redusă la verificări izolate ale unor metode individuale. Este necesară o abordare care să acopere atât validarea locală a regulilor de business, cât și comportamentul compus al aplicației în scenarii apropiate de utilizarea reală.

Capitolul de față este organizat în jurul infrastructurii de testare existente în proiect și al observațiilor rezultate din rularea și analiza componentelor implementate. Accentul principal este pus pe backend, unde există cea mai consistentă suită de teste automate, dar sunt discutate și gradul redus de acoperire din frontend, validările infrastructurale asociate serviciului Python și limitele actuale ale procesului de testare. În acest fel, evaluarea nu urmărește doar să confirme funcționarea unor scenarii fericite, ci și să identifice realist nivelul de maturitate tehnică atins de soluție.

\section{Strategia de testare și infrastructura utilizată}

În forma actuală a proiectului, strategia de testare este centrată pe backend-ul .NET, unde a fost construită o infrastructură clară pentru teste unitare, teste de controller și teste de integrare. Proiectul dedicat testelor utilizează framework-ul xUnit și include dependențe pentru mocking, pentru rularea aplicației în regim de test și pentru colectarea de coverage. Din această configurație rezultă o abordare în care pot fi validate atât componente individuale, cât și fluxuri complete la nivel HTTP, fără a depinde de o instanță externă a aplicației sau de o bază de date de producție.

Un rol central în această infrastructură îl are fabrica de aplicație folosită pentru testele de integrare. Aceasta pornește API-ul într-un mediu de tip \emph{IntegrationTesting}, injectează configurări specifice pentru conexiunea la bază de date, pentru JWT și pentru integrarea cu serviciul de recunoaștere, iar înlocuirea serviciului de email cu o implementare neutră permite verificarea fluxurilor principale fără efecte externe nedorite. Folosirea unei baze de date \emph{InMemory} pentru scenariile de integrare oferă izolare între teste și reduce costul de execuție, făcând posibilă validarea repetată a unor fluxuri complete precum înregistrarea, autentificarea sau aprobarea aplicațiilor de manager.

Din punct de vedere al contribuției tehnice, această infrastructură este importantă deoarece permite verificarea aplicației la mai multe niveluri fără a dubla artificial codul de test. În locul unei colecții disparate de verificări locale, proiectul beneficiază de o structură de testare coerentă, în care aceeași aplicație poate fi exercitată atât prin servicii și controllere izolate, cât și prin cereri HTTP apropiate de comportamentul real.

\section{Testarea unitară și validarea datelor}

Un prim nivel de validare este oferit de testele unitare orientate spre DTO-uri și spre reguli de validare a datelor de intrare. Aceste teste verifică faptul că obiectele folosite în autentificare, rezervări, recenzii, locuri de pescuit sau pontoane respectă constrângerile declarate prin atribute de validare. Din punct de vedere practic, această zonă este importantă deoarece multe dintre regulile de business pornesc de la date corecte și complete, iar acceptarea unor obiecte invalide ar produce erori în etapele ulterioare ale fluxurilor.

În implementarea actuală, aceste teste folosesc scenarii parametrizate pentru a acoperi multiple combinații de date valide și invalide. Sunt verificate, de exemplu, lungimea și forma credențialelor de autentificare, coerența dintre parolă și confirmarea acesteia, limitele pentru durata unei rezervări, intervalul permis pentru rating sau corectitudinea coordonatelor și a prețului unui loc de pescuit. Avantajul acestei abordări este că un număr relativ redus de metode de test poate acoperi un spectru larg de cazuri-limită, fără a fragmenta inutil suita de testare.

Pe lângă validarea DTO-urilor, testele unitare vizează și servicii cu logică de business mai densă, precum autentificarea și integrarea cu serviciul de analiză video. Aceste teste urmăresc comportamentul local al componentelor, izolat de infrastructura HTTP sau de persistență, și sunt utile pentru verificarea unor decizii precum generarea răspunsului de autentificare, tratarea erorilor sau actualizarea stărilor analizelor media. Din perspectiva evaluării, ele confirmă faptul că regulile esențiale sunt verificabile și în afara fluxurilor complete de integrare.

\section{Testarea controllerelor și a comportamentului API}

Un al doilea nivel de validare este oferit de testele dedicate controllerelor. Acestea verifică în mod direct răspunsurile HTTP, codurile de stare și reacția API-ului la cereri valide sau invalide, dar fără a porni întregul sistem în forma completă a testelor de integrare. Rolul lor este de a confirma că punctele de intrare ale backend-ului respectă contractele așteptate și că tratează corect atât situațiile reușite, cât și scenariile de eroare.

Pentru aplicația WhereWeFishin, această zonă este relevantă mai ales pentru controllerele care concentrează logică operațională, precum cele pentru autentificare, rezervări, administrarea locațiilor sau analiză video. În cazul rezervărilor, de exemplu, testarea controllerului permite verificarea modului în care sunt interpretate datele de intrare, a reacției la conflicte de disponibilitate și a integrării corecte cu componente auxiliare, precum notificările. În mod similar, controllerele de analiză video pot fi validate din perspectiva accesului autorizat, a limitelor de upload și a structurii răspunsurilor trimise către client.

Din punct de vedere metodologic, existența acestui nivel intermediar de testare este importantă deoarece reduce riscul ca problemele de contract API să fie descoperite abia în testele de integrare sau direct în utilizarea manuală. În plus, el oferă un cadru potrivit pentru verificarea controlată a claims-urilor utilizatorului și a regulilor de acces fără a fi nevoie de parcurgerea completă a tuturor fluxurilor de autentificare.

\section{Testarea de integrare a fluxurilor critice}

Zona cu cea mai mare valoare pentru evaluarea aplicației este reprezentată de testele de integrare. Acestea folosesc infrastructura comună a fabricii de aplicație și exercită API-ul prin cereri HTTP reale, validând atât răspunsurile expuse de backend, cât și efectele produse asupra datelor persistente. Din această perspectivă, testele de integrare oferă cea mai apropiată imagine de modul în care aplicația se comportă în exploatare normală.

Un prim exemplu relevant este fluxul complet de autentificare. Testele de integrare dedicate acoperă scenariul de înregistrare a unui utilizator nou, validarea persistenței sale, autentificarea ulterioară și verificarea tokenului obținut. În plus, aceeași suită tratează cazuri negative precum duplicarea username-ului sau a adresei de email, precum și transmiterea unor date invalide. Astfel, evaluarea nu se rezumă la verificarea traseului ideal, ci confirmă și robustețea mecanismelor de protecție împotriva datelor inconsistente.

Un al doilea exemplu important este workflow-ul de aplicare pentru manager. Acesta implică utilizatori cu roluri diferite și o succesiune de stări care trebuie menținute coerent. Testele de integrare dedicate acestei zone sunt valoroase deoarece validează nu doar structura răspunsurilor API, ci și corectitudinea tranzițiilor de stare și impactul deciziilor administrative asupra datelor persistente. Din punct de vedere al capitolului de evaluare, acest flux este una dintre cele mai bune dovezi că partea de business logic implementată în backend este verificată pe scenarii relevante și nu doar pe fragmente izolate de cod.

Pe aceeași linie se află și fluxurile legate de rezervări și roluri operaționale. Chiar dacă nu toate scenariile sunt acoperite prin aceeași adâncime de integrare, prezența testelor pentru booking și pentru module precum angajații sau administrarea platformei arată că sistemul este evaluat și în situațiile în care mai multe componente trebuie să colaboreze: repository-uri, servicii, autorizare și baze de date de test. Această abordare crește încrederea în faptul că aplicația răspunde coerent atunci când regulile de business sunt exercitate cap-coadă.

\section{Validarea operațională a comportamentului sistemului}

Dincolo de testele automate, evaluarea aplicației trebuie să țină cont și de mecanismele de protecție și reziliență configurate direct în infrastructura runtime. În cazul backend-ului WhereWeFishin, aceste elemente sunt vizibile în configurația principală a aplicației și oferă o formă de validare operațională complementară testării clasice.

Un prim element important este politica strictă asociată JWT-urilor. Validarea explicită a issuer-ului, a audience-ului, a semnăturii și a duratei de viață, împreună cu o fereastră de toleranță nulă pentru expirare, contribuie la un model de securitate predictibil. În paralel, endpoint-urile sensibile de autentificare sunt protejate prin rate limiting, ceea ce reduce riscul de abuz și arată că aplicația a fost gândită și din perspectivă operațională, nu doar funcțională.

Un al doilea element relevant este legat de manipularea cererilor voluminoase și a fluxurilor media. Limitarea dimensiunii cererilor la 150 MB, împreună cu configurarea explicită a timeout-urilor și a conexiunilor, oferă un cadru de exploatare mai stabil pentru încărcarea de videoclipuri. Tot aici pot fi incluse compresia răspunsurilor, politicile de output caching și mecanismele de retry la nivelul conexiunii cu baza de date, toate contribuind la o aplicație mai robustă în condiții reale de utilizare.

Această evaluare operațională nu înlocuiește testele automate, dar completează imaginea de ansamblu. Ea arată că soluția nu a fost construită doar pentru a trece teste locale, ci și pentru a răspunde mai bine unor cerințe reale privind securitatea, traficul și reziliența infrastructurii.

\section{Acoperire actuală, limitări și direcții de extindere}

O evaluare realistă trebuie să evidențieze nu doar punctele forte ale procesului de testare, ci și limitele actuale ale acestuia. În proiectul analizat, backend-ul beneficiază de cea mai bună acoperire, datorită combinației dintre teste unitare, teste de controller și teste de integrare. În schimb, partea de frontend este reprezentată printr-o acoperire automată redusă, centrată în principal pe un serviciu specific, iar serviciul Python pentru analiză media dispune mai ales de verificări infrastructurale și de tip smoke test, nu de o suită bogată de teste funcționale automate.

Această distribuție inegală a acoperirii este de înțeles într-un proiect în care complexitatea logicii de business este concentrată în backend, însă reprezintă totuși o limitare importantă. Lipsa unei suite consistente de teste pentru interfață înseamnă că multe dintre fluxurile frontend sunt validate în practică mai ales prin integrarea cu backend-ul și prin verificări manuale. În mod similar, partea de analiză media rămâne dependentă într-o măsură mai mare de validarea experimentală și de observarea comportamentului în execuție, decât de teste unitare fine pentru fiecare etapă din pipeline.

O altă limitare este absența, în forma actuală, a unor benchmark-uri explicite de performanță sau a unor teste end-to-end complete care să includă simultan interfața, backend-ul și serviciile externe. Din acest motiv, capitolul de față poate argumenta convingător corectitudinea funcțională și robustețea arhitecturală a unor module importante, dar nu poate pretinde o evaluare exhaustivă din perspectiva scalării sau a experienței utilizatorului în condiții de sarcină ridicată.

Ca direcții de extindere, pot fi avute în vedere dezvoltarea unei suite mai consistente de teste frontend, introducerea unor teste automatizate pentru pipeline-ul Python și generarea sistematică a unor rapoarte de coverage care să poată fi discutate și cantitativ. Chiar și în lipsa acestor extinderi, infrastructura existentă oferă o bază solidă pentru validarea logicii esențiale a platformei.

\section{Sinteza capitolului}

Testarea și evaluarea soluției WhereWeFishin arată că aplicația beneficiază de o bază solidă de validare în special la nivel de backend, unde sunt acoperite atât regulile de validare a datelor, cât și comportamentul controllerelor și fluxurile de integrare cele mai importante. Infrastructura de test construită în jurul xUnit și a fabricii de aplicație permite verificarea unor scenarii realiste, precum autentificarea completă, workflow-ul de manager application și colaborarea dintre mai multe componente ale sistemului.

În același timp, evaluarea evidențiază limitele curente ale proiectului: acoperire automată redusă în frontend, validare limitată a serviciului Python și absența unor benchmark-uri de performanță mai ample. Aceste observații nu diminuează valoarea implementării, ci oferă o imagine realistă asupra maturității tehnice atinse și asupra pașilor care pot urma. Pe această bază, capitolul final poate sintetiza contribuțiile proiectului și direcțiile sale de dezvoltare ulterioară.